\section{RESULTS}

This section is intentionally structured before values are inserted so that
result selection follows the hypotheses rather than whichever demonstrations
look best.

\subsection{Flat-Ground Transport}

\begin{table}[t]
  \caption{Flat-ground performance on the same platform}
  \label{tab:flat-results}
  \centering
  \small
  \begin{tabular}{lcccc}
    \toprule
    Mode & $N$ & Speed & Success & Energy/m \\
    \midrule
    Articulated walk & TBD & TBD & TBD & TBD \\
    Distal-only & TBD & TBD & TBD & TBD \\
    Full quasi-roll & TBD & TBD & TBD & TBD \\
    \bottomrule
  \end{tabular}
\end{table}

\draftnote{Report means with standard deviations and units in every header. The
2.0 m/s and 0.7 m/s historical claims enter the paper only after raw measurement
reconstruction or a repeated controlled test.}

\subsection{Stair Traversal}

\begin{figure}[t]
  \centering
  \fbox{\parbox[c][32mm][c]{0.92\columnwidth}{\centering
    Final Fig.~4: success rate versus normalized riser height, plus ascent time
    for the 15-step staircase. Show raw trial points and uncertainty.}}
  \caption{Stair-climbing performance under onboard power.}
  \label{fig:stair-results}
\end{figure}

\subsection{Ablation and Failure Analysis}

The ablation must answer whether articulated body/contact adaptation improves
stair performance beyond rotating arcs alone. Failure categories should include
slip, missed edge engagement, actuator overload, body collision, loss of balance,
and controller or power faults. Representative failures belong in the submitted
video as well as the success examples.

\begin{figure}[t]
  \centering
  \fbox{\parbox[c][28mm][c]{0.92\columnwidth}{\centering
    Final Fig.~5: current, body pitch, and distal phase over one successful and one
    failed ascent.}}
  \caption{Time-aligned state and electrical measurements reveal how failures
  arise rather than merely reporting whether a run succeeded.}
  \label{fig:failure-results}
\end{figure}
